% ATTEMPT_STATUS: unresolved
% TOP_PROBLEM: {"problemId": "problem.top500-omission-w041044-weight-monodromy-conjecture", "canonicalTitle": "Weight–Monodromy Conjecture", "releaseRank": 85, "primaryDomain": "number_theory_arithmetic_geometry", "primaryDomainLabel": "Number theory & arithmetic geometry", "displayStatus": "open_with_solved_subcases", "releaseStatus": "open_with_solved_subcases"}
% !TeX program = lualatex
\documentclass[11pt]{article}
\usepackage[margin=25mm]{geometry}
\usepackage{fontspec}
\setmainfont{Latin Modern Roman}
\usepackage{amsmath,amssymb,mathtools}
\usepackage{unicode-math}
\setmathfont{Latin Modern Math}
\usepackage{xurl,hyperref}
\hypersetup{colorlinks=true,urlcolor=blue}
\setcounter{secnumdepth}{0}
\setlength{\parindent}{0pt}
\setlength{\parskip}{5pt}
\setlength{\emergencystretch}{3em}
\begin{document}
\section*{\texorpdfstring{Weight–Monodromy Conjecture}{Weight–Monodromy Conjecture}}\label{85}
\subsection{Definitions and mathematical statement}
Put $V=H^i_{\mathrm{\acute et}}(X_{\bar K},\overline{\mathbb{Q}}_\ell)$ for a proper smooth variety over a nonarchimedean local field with residue cardinality $q$, and $\ell\ne p$. Inertia defines nilpotent monodromy $N:V\to V(-1)$. Its zero-centered increasing filtration $M$ is characterized by
\[
 N(M_j)\subseteq M_{j-2}(-1),\qquad
 N^a:\operatorname{gr}^M_aV\xrightarrow{\sim}\operatorname{gr}^M_{-a}V(-a).
\]
For any geometric Frobenius lift, every eigenvalue $\alpha$ on $\operatorname{gr}^M_jV$ should be algebraic and satisfy $|\iota(\alpha)|=q^{(i+j)/2}$ for every complex embedding $\iota$. This is purity of weight $i+j$. The original field and its $q$ remain fixed when defining monodromy on an open inertia subgroup; no semistability or projectivity assumption is added.
\par\smallskip\textit{Formulation and concept references:} \href{https://arxiv.org/pdf/1111.4914v1}{[S1]}.

\subsection{Short English statement}
Do monodromy layers of local-field étale cohomology have exactly the Frobenius weights predicted by their degrees?
\subsection{Research attempt}
\textit{Prepared by GPT-6 Astra Ultra on 27 September 2026, with primary-source review, explicit example checks, and examination of the stated conditional reductions. This is an unresolved research attempt.}

\paragraph{Scope and source check (27 September 2026).}
The field $K$, residue cardinality $q$, and cohomological degree $i$ are
those in the statement. No projectivity or semistable model is assumed
for the target. Scholze [S1, Conjecture 1.13 and Theorem 1.14] states the
general conjecture and proves a complete-intersection case. Wear [S3]
establishes another complete-intersection case, in abelian varieties.
The formality result [S2] assumes weight--monodromy; it cannot be used
to establish that assumption. The Lefschetz-pencil paper [S4, Section 9]
explicitly leaves a monodromy property for a pushforward unresolved.
These are relevant results and reductions, not a proof in the full scope.

This attempt separates the formal nilpotent linear algebra from the
geometric purity statement. It proves a reduction to primitive monodromy
pieces, checks the Frobenius and twist signs, and explains precisely why
passing to a finite extension preserves the claim. The remaining problem
will be the weights of the primitive pieces, which are not determined
by the Jordan form of monodromy.

\paragraph{The sign convention in one Jordan chain.}
Choose a basis for the Tate line when writing matrices. Geometric
Frobenius $F$ acts on $\mathbb Q_\ell(1)$ by $q^{-1}$, so it acts
on the twist $V(-1)$ by $qF$. Equivariance of
$N:V\to V(-1)$ therefore means
\begin{equation}\label{wm:relation}
 NF=qFN,\qquad FNF^{-1}=q^{-1}N.
\end{equation}
The untwisted underlying operator lowers Frobenius eigenvalues by a
factor of $q$. Reversing this factor would reverse the predicted weights.

For a chain $e_0,\ldots,e_r$ with $Ne_a=e_{a-1}$ and $Ne_0=0$,
the zero-centered monodromy grades are
\[
 \deg_M e_a=-r+2a.
\]
One compatible diagonal Frobenius is $Fe_a=\lambda q^a e_a$.
The conjecture on this chain is exactly
\[
 |\iota(\lambda)|=q^{(i-r)/2};
\]
then $\lambda q^a$ has weight $i-r+2a$, as required. For a length-two
chain in degree one, the familiar eigenvalues $1,q$ belong to grades
$-1,1$. This is also a direct test of the displayed isomorphism
$N:\operatorname{gr}_1\to\operatorname{gr}_{-1}(-1)$: both sides have
Frobenius eigenvalue $q$. A twist by $+1$ in that isomorphism would fail.

Nothing here assumes that all Frobenius operators are diagonalizable.
Equation \eqref{wm:relation} maps a generalized $\alpha$-eigenspace
to the generalized $\alpha/q$-eigenspace. Statements about eigenvalue
weights must not be strengthened to semisimplicity without a separate
argument.

\paragraph{A complete reduction to the primitive graded pieces.}
For $r\ge0$, put
\[
 P_r=\ker\left(N^{r+1}:\operatorname{gr}^M_rV
                \longrightarrow\operatorname{gr}^M_{-r-2}V(-r-1)\right).
\]
These are Frobenius-stable spaces. Nilpotent linear algebra gives
\begin{equation}\label{wm:primitive}
 \operatorname{gr}^M_j V\simeq
 \bigoplus_{\substack{r\ge |j|\\r\equiv j\pmod2}}
 P_r\bigl((r-j)/2\bigr).
\end{equation}
Here the map from a summand is the appropriate power of $N$, interpreted
with its Tate twist. To verify the statement without assuming a
Frobenius-stable Jordan basis, first choose an ordinary Jordan basis for
$N$. A chain of length $r+1$ contributes once in grades
$r,r-2,\ldots,-r$. In grade $r$, the kernel defining $P_r$ removes
the contributions from every longer chain: on those contributions
$N^{r+1}$ is injective. On a chain of length $r+1$ it vanishes.
Thus the indicated maps give a direct sum decomposition in every grade.
The maps and kernels themselves are intrinsic and Frobenius-equivariant,
so the resulting decomposition is one of Frobenius representations.

If $P_r$ is pure of weight $i+r$, its twist in
\eqref{wm:primitive} is pure of weight
\[
 i+r-2\frac{r-j}{2}=i+j.
\]
Conversely, if all the graded pieces have the conjectured weights,
their subspaces $P_r$ do also. Thus the conjecture is equivalent to
purity of these primitive spaces. This reduction does not construct
their eigenvalues or prove they are algebraic; both requirements remain
part of the purity assertion.

The identity also gives exact dimension checks. A chain of length
$r+1$ contributes one dimension to $P_r$, so
\[
 \dim V=\sum_{r\ge0}(r+1)\dim P_r,\qquad
 \dim\operatorname{gr}_j^M V=\dim\operatorname{gr}_{-j}^M V.
\]
These numerical symmetries are consequences of monodromy alone. They
are necessary but do not prove a statement about complex absolute values.

\paragraph{Changing a Frobenius lift does not change the purity test.}
An open inertia subgroup acts as $\exp(t_\ell(\sigma)N)$.
Since $N$ strictly lowers the monodromy filtration, this subgroup acts
trivially on its associated graded. Therefore the residual inertia
image $H$ on the total graded space is finite. A Frobenius operator
normalizes $H$.

Suppose $F'=hF$ is another geometric lift, with $h\in H$. Choose
$a>0$ so that conjugation by $F^a$ is trivial on $H$; this is possible
because $\operatorname{Aut}(H)$ is finite. Then
$(hF)^a=h_aF^a$ for some $h_a\in H$, and $h_a$ commutes with $F^a$.
If $b$ is an exponent of the finite group $H$, it follows that
\[
 (F')^{ab}=F^{ab}
\]
on the associated graded. Thus the multisets of eigenvalues of the
two operators have the same $ab$-th powers. Algebraicity and every
specified absolute-value weight for one multiset imply them for the
other. Indeed an element whose positive integral power is algebraic
is algebraic, and taking absolute values recovers the unique positive
$ab$-th root. This proves the lift independence needed here without
pretending that inertia was already trivial on $V$.

\paragraph{Finite extensions and the original residue cardinality.}
Let $K'/K$ have residue degree $f$, so $q'=q^f$. Restriction rescales
the underlying monodromy operator by a nonzero scalar if a different
tame generator is used; kernels and images of its powers, and hence
the zero-centered filtration, do not change. A geometric Frobenius
for $K'$ acts on the graded space as $hF^f$ for a finite inertia
element $h$. The preceding argument, applied to $F^f$, compares their
positive powers.

Consequently purity over $K'$ of weight $i+j$ is equivalent to purity
over $K$ of the same weight. On eigenvalues the numerical calculation is
\[
 |\iota(\alpha^f)|=(q^f)^{(i+j)/2}
 \quad\Longleftrightarrow\quad
 |\iota(\alpha)|=q^{(i+j)/2}.
\]
Algebraicity descends in the same way, and every embedding of the field
generated by $\alpha$ restricts to an embedding of the field generated
by $\alpha^f$. The additional finite-inertia powers cause no change.
This justifies an extension used in a geometric argument; it does not
allow changing $q$ while keeping the same Frobenius eigenvalue.

\paragraph{A genuine solved case and an algebraic countertest.}
If $X$ has a proper smooth model over the valuation ring, smooth proper
base change identifies $V$ with the cohomology of the special fiber.
Inertia is trivial, $N=0$, and only $\operatorname{gr}_0^M V$ survives.
The finite-field purity theorem gives precisely weight $i$. These
established geometric inputs, recalled in [S1], prove the good-reduction
case. Merely knowing $N=0$ for an abstract representation would not
give finite-field geometric purity.

For a concrete test, take two length-two chains in formal degree $i=1$,
and set
\[
 F=\operatorname{diag}(q^s,q^{s+1},q^{-s},q^{1-s}),
 \qquad Ne_1=e_0,\quad Nf_1=f_0,
\]
in the basis $(e_0,e_1,f_0,f_1)$, with $s=1$, for instance.
Equation \eqref{wm:relation} holds, both monodromy grades have
dimension two, and $\det F=q^2$, the determinant size expected for
four dimensions of weight one on average. Yet grade $-1$ has eigenvalues
$q^s,q^{-s}$, which do not both have weight zero. The top grade fails
in the corresponding way. This is a formal Frobenius--monodromy pair,
not a claimed cohomology representation or a counterexample to the
conjecture. It shows exactly why nilpotence, the chain isomorphisms,
and even the total determinant do not imply purity.

\paragraph{Verification and the geometric gap.}
The local checks construct chains of several lengths and verify
$FNF^{-1}=q^{-1}N$, the grade dimensions, the primitive count, and the
four-dimensional determinant countertest with rational matrices.
They also check the exponents after a residue-degree extension. These
calculations audit signs and linear algebra, not geometric realizability.

The first missing statement is that the actual geometric $P_r$ are
pure of weight $i+r$. In a semistable spectral-sequence approach,
purity of individual special-fiber strata is not enough: one must
control the induced monodromy maps after taking cohomology of the
differentials. The unresolved pushforward condition in [S4] is a
precise manifestation of that issue. Possible continuation would
prove the needed pairing nondegeneracy in a specified geometric
family, produce a comparison preserving these primitive pieces, or
establish the monodromy isomorphisms on the relevant spectral-sequence
page. The linear algebra here supplies no such theorem.

\textbf{UNRESOLVED.} The attempt proves formal reductions and normalization
checks, and explains an established good-reduction subcase. It neither
adds a hypothesis to the catalogue target nor proves its remaining
mixed-characteristic cases. No novelty or formal proof certification
is asserted.

\subsection{Sources}
\begin{itemize}
\item[S1]Peter Scholze, Perfectoid spaces, arXiv1111.4914v1,banner21November2011/internal27November2024,51pages.
\url{https://arxiv.org/pdf/1111.4914v1}.
\textit{Formulation source.}
\item[S2]Formality for rigid-analytic spaces satisfying the weight-monodromy conjecture.
\url{https://arxiv.org/abs/2607.14517}.
\textit{Further reference; not a proof endorsement.}
\item[S3]Perfectoid covers of abelian varieties and the weight-monodromy conjecture.
\url{https://arxiv.org/abs/2303.05610}.
\textit{Further reference; not a proof endorsement.}
\item[S4]Semistable Lefschetz pencils.
\url{https://page.mi.fu-berlin.de/esnault/preprints/helene/151_mw.pdf}.
\textit{Further reference; not a proof endorsement.}
\end{itemize}
\end{document}
